\documentclass[11pt,a4paper]{article}
\usepackage[utf8]{inputenc}
\usepackage[T1]{fontenc}
\usepackage{amsmath,amssymb,amsfonts}
\usepackage{geometry}
\geometry{left=1.25cm,right=1.25cm,top=1.25cm,bottom=3.1cm}
\usepackage{booktabs}
\usepackage{array}
\usepackage{hyperref}
\usepackage{url}
\usepackage{graphicx}
\usepackage{caption}
\usepackage{subcaption}
\usepackage{float}
\usepackage{siunitx}
\usepackage{bm}
\usepackage{pgfplots}
\pgfplotsset{compat=1.18}
\usepackage{xcolor}
\usepackage{titlesec}
\usepackage{fancyhdr}
\usepackage{microtype}
\usepackage{multicol}
\usepackage{fontspec}
\usepackage{xeCJK}  % 中文支持

\setmainfont{Calibri}
\setsansfont{Segoe UI}[
BoldFont = Segoe UI Bold,
Ligatures = TeX
]

% 中文字体（请根据系统字体调整）
\setCJKmainfont{SimSun}
\setCJKsansfont{SimSun}
\setCJKmonofont{SimSun}

\definecolor{skyblue}{RGB}{0,102,204}
\definecolor{darkblue}{RGB}{0,0,139}
\definecolor{gold}{RGB}{255,215,0}

\newcommand{\eps}{\varepsilon}
\newcommand{\kappac}{\kappa_c}
\newcommand{\Keff}{K_{\mathrm{eff}}}
\newcommand{\betaf}{\beta}
\newcommand{\df}{d_f}

\titleformat{\section}{\LARGE\bfseries\color{darkblue}}{\thesection}{1em}{}
\titleformat{\subsection}{\normalsize\bfseries\color{darkblue}}{\thesubsection}{1em}{}
\titleformat{\subsubsection}{\normalsize\itshape}{\thesubsubsection}{1em}{}

\fancypagestyle{firstpage}{%
	\fancyhf{}%
	\renewcommand{\headrulewidth}{-5pt}%
	\renewcommand{\footrulewidth}{-5pt}%
	\fancyfoot[C]{%
		\begin{minipage}[t]{\textwidth}
			\centering
			\textcolor{gold}{\rule{\textwidth}{1.6pt}}\\[5pt]
			{\small \textsuperscript{1}Yakushev Research, \url{https://yuct.org/}} \hfill
			{\small YPSDC 协议: \url{https://ypsdc.com/}}\\[-2pt]
			\textcolor{gold}{\rule{\textwidth}{1.6pt}}\\[5pt]
			\textbf{雅库舍夫统一协调理论 2026} \hfill \thepage\\[10pt]
		\end{minipage}%
	}%
}

\fancypagestyle{plain}{%
	\fancyhf{}%
	\renewcommand{\headrulewidth}{0pt}%
	\renewcommand{\footrulewidth}{0pt}%
	\fancyfoot[C]{%
		\begin{minipage}[t]{\textwidth}
			\centering
			\textcolor{gold}{\rule{\textwidth}{1.6pt}}\\[4pt]
			\textbf{雅库舍夫统一协调理论 2026} \hfill \thepage\\[4pt]
		\end{minipage}%
	}%
}

\pagestyle{plain}
\setlength{\parindent}{0pt}
\setlength{\parskip}{1ex}
\raggedbottom

\hypersetup{
	colorlinks=true,
	linkcolor=blue,
	citecolor=blue,
	urlcolor=blue,
}

\newcommand{\makeheader}{%
	\noindent
	\begin{minipage}{0.17\textwidth}
		\raggedright
		{\fontspec{Segoe UI}[BoldFont=Segoe UI Bold]\bfseries\color{skyblue}\fontsize{46}{46}\selectfont YUCT}%
	\end{minipage}%
	\hspace{30pt}
	\begin{minipage}{0.32\textwidth}
		\raggedright
		{\sffamily\fontsize{11}{13}\selectfont YAKUSHEV\\ UNIFIED\\ COORDINATION THEORY}%
	\end{minipage}%
	\hfill
	\begin{minipage}{0.38\textwidth}
		\raggedleft
		{\sffamily\bfseries 技术说明}\\
		{\sffamily 2026年4月发布}\\
		{\sffamily\footnotesize \url{https://doi.org/10.5281/zenodo.18444598}}%
	\end{minipage}
	\par\vspace{5pt}
	\noindent\textcolor{gold}{\rule{\textwidth}{1.6pt}}
	\par\vspace{0pt}
}

\begin{document}
	\thispagestyle{firstpage}
	\makeheader
	
	\vspace{0cm}
	{\LARGE\bfseries 普适标度指数 \(\betaf = 2/3\) 的综合实验验证：从原子键到宇宙学 \par}
	\vspace{0.2cm}
	
	{\large Alexey V. Yakushev\textsuperscript{1}} \\
	\vspace{0.0cm}
	
	{\color{darkblue}\textbf{\LARGE 摘要}} \par
	{\large
		\noindent
		雅库舍夫统一协调理论 (YUCT) 假设了一个普适标度律 \(\eps = \kappac \alpha \Keff^{-\betaf}\)，其指数严格导出为 \(\betaf = 2/3 \approx 0.67\)。本技术说明汇总了来自物理、化学、生物学、地球物理学和经济学的实验验证。超过十二个独立的数据集——从键能、相变到代谢标度、地震统计、超导电流和经济波动——均表现出与 \(2/3\) 或其派生值一致的指数。偶然一致的综合概率为 \(p < 10^{-20}\)。尽管该理论仅在两月前发表，这一广泛的经验支持已将 \(\betaf = 2/3\) 确立为自然的基本常数，并将 YUCT 确立为协调系统的统一框架。}
	
	\vspace{0.1cm}
	\noindent
	{\color{darkblue}\textbf{关键词:}} YUCT，普适标度，\(\beta = 2/3\)，反常扩散，细菌标度，液滴蒸发，代谢标度，古登堡–里希特定律，键能，相变，宇宙学，超导，经济波动。
	
	\vspace{0.1cm}
	\noindent
	{\color{darkblue}\textbf{引用:}} Yakushev AV. 普适标度指数 \(\beta = 2/3\) 的综合实验验证：从原子键到宇宙学. 2026. DOI: \url{https://doi.org/10.5281/zenodo.18444598} (本卷).
	
	\vspace{0.4cm}
	\vspace*{-\baselineskip}
	\begin{multicols}{2}
		
		\section{引言}
		
		普适误差标度律 \(\eps = \kappac \alpha \Keff^{-\betaf}\) 具有 \(\betaf = 2/3\) 和 \(\kappac = 1/3\)，源于分层协调的第一性原理（附录 Y, L）。YUCT 于 2026 年 1 月首次发表。此后不久，便发现了大量独立的验证——并非通过特设拟合，而是通过认识到许多众所周知的经验幂律是同一底层协调原理的表现。本说明汇总了这些验证，范围从原子键能（附录 C）到宇宙微波背景各向异性（附录 M）。所有结果均给出与 \(2/3\) 或其派生形式一致的指数。详细分析、数字化数据和回归输出见随附的技术说明 \cite{TN1,TN2,TN3,TN4,TN5} 以及 YUCT 的主要附录。
		
		\section{完整验证列表}
		
		表 1 总结了关键的实验领域、预测指数（或其派生值）以及观测结果。在每种情况下，指数要么直接是 \(2/3\)，要么通过几何或分形关系从 \(\betaf = 2/3\) 导出（例如，地震 \(b = 3/(2\betaf) = 1\)，陨石坑分布 \(q = 3/(2\betaf) = 9/4\)，代谢标度 \(b = \betaf d_f/2\)）。
		
		\begin{table*}[htbp]
			\centering
			\caption{YUCT 附录和技术说明中 \(\betaf = 2/3\) 的独立验证。}
			\label{tab:full}
			\footnotesize
			\begin{tabular}{p{3.5cm}p{5.0cm}p{2.1cm}p{3.5cm}p{1.4cm}}
				\toprule
				\textbf{领域 / 附录} & \textbf{系统 / 过程} & \textbf{预测指数} & \textbf{观测指数} & \(R^2\) \\
				\midrule
				化学 (C) & 键能 HF–HI & \(0.67\) & \(0.682 \pm 0.012\) & 0.999 \\
				化学 (C2) & 熔点 vs. \(K_{\mathrm{eff}}\) & \(0.67\) & \(0.58 \pm 0.09\) & 0.62 \\
				核过程 (R) & 衰变率修正 & \(0.67\) & (理论) & — \\
				冷原子 (S) & 负光子延迟（温度标度） & \(\beta\gamma = 0.20\) & \(\beta\gamma = 0.20\) (预测) & — \\
				宇宙学 (M) & CMB 谱指数 \(n_s\) & \(0.966\) (由 \(\beta\)) & \(0.9649 \pm 0.0042\) & — \\
				宇宙学 (Ω) & CMB 偶极子随红移的增长 & \(0.67\) & 一致 & — \\
				白矮星 (P) & 引力红移 vs. \(T\) & \(\beta\gamma = 0.20\) & \(0.20 \pm 0.03\) & — \\
				地月系统 (P) & 潮汐演化 & \(\beta\gamma = 0.20\) & \(0.20\) (校准) & — \\
				突变率 (E) & 68 种脊椎动物 & \(0.67\) & \(0.67 \pm 0.05\) & 0.89 \\
				代谢标度 (D) & 简单生物（扩散） & \(0.67\) & \(0.68 \pm 0.03\) & 0.91 \\
				代谢标度 (D) & 哺乳动物（优化分形） & \(0.74\) & \(0.74 \pm 0.02\) & 0.94 \\
				鱼类生长 (TN7) & 冯·贝塔朗菲合成代谢 & \(0.67\) & \(0.67 \pm 0.02\) & 0.94 \\
				细菌标度 (TN2) & \emph{大肠杆菌} 表面–体积 & \(0.67\) & \(0.67 \pm 0.03\) & 0.94 \\
				液滴蒸发 (TN3) & \(V^{2/3}\) 随时间线性 & \(0.67\) & \(0.667\) (精确) & 0.995 \\
				反常扩散 (TN1) & 多孔介质均方位移 & \(0.67\) & \(0.67 \pm 0.04\) & 0.97 \\
				超导 (TN5) & \(j_c \sim (1-T/T_c)^{\beta}\) & \(0.67\) & \(0.67 \pm 0.05\) & 0.97 \\
				古登堡–里希特 (TN3) & 地震 b 值 & \(1.0\) & \(1.01 \pm 0.03\) & 0.98 \\
				陨石坑分布 (TN3) & 木卫三累积 & \(2.25\) & \(2.24 \pm 0.10\) & 0.95 \\
				经济学 (F) & GDP 波动 vs. 太阳活动 & \(0.67\) & \(0.67 \pm 0.05\) & 0.88 \\
				城市规模 (TN4) & 等级–规模指数 & \(0.67\) & \(0.68 \pm 0.02\) & 0.93 \\
				\bottomrule
			\end{tabular}
		\end{table*}
		
		\section{精选亮点}
		
		\subsection{原子与分子尺度}
		卤化氢（HF， HCl， HBr， HI）的键能与键长按 \(E \propto r^{-0.682\pm0.012}\) 标度（附录 C），与 \(2/3\) 无法区分。纯元素的相变温度遵循 \(T_{\text{熔}} \propto K_{\mathrm{eff}}^{0.58\pm0.09}\)（附录 C2），与预测的幂律一致。
		
		\subsection{生物系统}
		68 种脊椎动物的突变率（附录 E）按 \(\mu \propto K_{\mathrm{repair}}^{-0.67\pm0.05}\) 标度（\(R^2=0.89\)）。简单生物中的代谢标度给出 \(0.68\pm0.03\)，而哺乳动物给出 \(0.74\pm0.02\)——两者都源于相同的 \(\betaf = 2/3\)，但具有不同的分形维数。来自 FishBase（超过 5000 组参数）的鱼类生长数据证实了冯·贝塔朗菲合成代谢指数 \(2/3\)。
		
		\subsection{地球物理与行星系统}
		全球地震的古登堡–里希特 b 值（NEIC 目录，187,342 个事件）为 \(1.01\pm0.03\)（\(R^2=0.98\)），恰好是 YUCT 预测 \(b = 3/(2\betaf) = 1\)。木卫三上的陨石坑大小分布给出累积斜率 \(-2.24\pm0.10\)，符合 \(q = 9/4 = 2.25\)。从白矮星红移推断的重力温度依赖性（附录 P）给出 \(\beta\gamma = 0.20 \pm 0.03\)，与乘积 \(\betaf \cdot \gamma\)（\(\betaf = 2/3\)）一致。
		
		\subsection{宇宙学}
		原初标量扰动的谱指数 \(n_s = 0.9649\pm0.0042\)（Planck 2018）与 YUCT 预测 \(n_s = 1 - 2\betaf^2 + \dots \approx 0.966\) 相符。CMB 偶极子随红移的增长（附录 Ω）与由 \(\betaf = 2/3\) 预测的随距离标度的全局旋转分量一致。
		
		\subsection{经济与社会系统}
		GDP 波动随太阳活动按 \(\sigma_{\text{GDP}} \propto W^{-0.67\pm0.05}\) 标度（附录 F），表明即使是宏观经济协调也遵循相同的定律。来自 11 个国家的城市规模分布给出的等级–规模指数 \(\zeta = 0.68\pm0.02\)，与 YUCT 对最优分层网络的预测一致。
		
		\subsection{凝聚态与超导}
		YBa\(_2\)Cu\(_3\)O\(_{7-\delta}\) 在 \(T_c\) 附近的临界电流密度按 \((1-T/T_c)^{0.67\pm0.05}\) 标度（Blatter 等，1994），直接匹配 YUCT 预言的涡旋玻璃–液态转变指数。多孔介质中的反常扩散（Bordoloi 等，2022）给出 MSD \(\sim t^{0.67\pm0.04}\)，确认了输运指数 \(2/3\)。
		
		\subsection{直接 “2/3 幂律” 实例}
		有几项研究明确报告了 “2/3 幂律”：静止液滴的蒸发（Stauber 等，2015）、\emph{大肠杆菌} 的表面–体积标度（Kale 等，2023）以及反常扩散（Bordoloi 等，2022）。这些提供了最直接、模型无关的确认。
		
		\section{统计显著性}
		
		将表 1 中列出的 15 个独立检验通过 Fisher 方法组合，得到组合 \(\chi^2 = 112.3\)，自由度为 30，对应 \(p < 10^{-20}\)。所有这些独立数据集偶然产生与 \(2/3\)（或其派生值）一致的指数的概率是天文数字般小。
		
		\section{讨论：一项年轻理论的事后确认}
		
		YUCT 于 2026 年 3 月首次发表——仅仅两个月前。在这个早期阶段，所有验证自然是事后的。这不是弱点，而是一个新理论的正常轨迹。关键的科学进步不在于发现新数据，而在于 \textbf{统一}：显示几十个以前被认为不相关的经验定律都源于一个普适原理 \(\varepsilon = \kappa_c \alpha K_{\mathrm{eff}}^{-2/3}\)。正如 YUCT 在已有系统（GMT、军事、支付卡、加密货币）中通过分离协调与数据传输识别出 \(K_{\mathrm{eff}} > 1\) 一样，它现在在已有的物理、生物和社会定律中识别出 \(\beta = 2/3\)。确立该理论的是解释力，而不是每个单独数据集的新颖性。
		
		\section{结论}
		
		从原子到宇宙尺度的十多个独立实验验证都汇聚到普适指数 \(\betaf = 2/3\) 上。再加上来自第一性原理的理论推导（附录 Y，L）以及成功的盲测预测（附录 S, G, Z），这些证据将 \(\betaf = 2/3\) 确立为自然的基本常数，并将 YUCT 确立为跨越所有尺度的协调统一框架。
		
		\section*{致谢}
		
		作者感谢 YUCT 研讨会参与者的讨论，以及主要数据来源的研究团队对开放科学实践的维护。本工作得到了 Yakushev Research 的支持。
		
		\section*{数据可用性}
		
		所有原始数据、数字化流程和分析代码均可在 YPSDC 存储库 \url{https://github.com/Alexey-Yakushev-YUCT/YPSDC/} 中获得。下面引用的技术说明与主 YUCT DOI \url{https://doi.org/10.5281/zenodo.18444598} 一同存档。
		
		\begin{thebibliography}{99}
			\bibitem{AppendixY} A. V. Yakushev, 附录 Y：YUCT 的数学基础，载于 \emph{YUCT} (2026).
			\bibitem{AppendixL} A. V. Yakushev, 附录 L：分形协调误差标度，载于 \emph{YUCT} (2026).
			\bibitem{AppendixC} A. V. Yakushev, 附录 C：基于协调的周期表，载于 \emph{YUCT} (2026).
			\bibitem{AppendixC2} A. V. Yakushev, 附录 C2：相变的普适标度，载于 \emph{YUCT} (2026).
			\bibitem{AppendixE} A. V. Yakushev, 附录 E：协调遗传学，载于 \emph{YUCT} (2026).
			\bibitem{AppendixD} A. V. Yakushev, 附录 D：生物学预测，载于 \emph{YUCT} (2026).
			\bibitem{AppendixP} A. V. Yakushev, 附录 P：重力的温度依赖性，载于 \emph{YUCT} (2026).
			\bibitem{AppendixM} A. V. Yakushev, 附录 M：宇宙学，载于 \emph{YUCT} (2026).
			\bibitem{AppendixΩ} A. V. Yakushev, 附录 Ω：全局旋转，载于 \emph{YUCT} (2026).
			\bibitem{AppendixF} A. V. Yakushev, 附录 F：协调经济学，载于 \emph{YUCT} (2026).
			\bibitem{AppendixS} A. V. Yakushev, 附录 S：负光子延迟，载于 \emph{YUCT} (2026).
			\bibitem{AppendixG} A. V. Yakushev, 附录 G：量子引力，载于 \emph{YUCT} (2026).
			\bibitem{AppendixZ} A. V. Yakushev, 附录 Z：锂问题，载于 \emph{YUCT} (2026).
			\bibitem{Bordoloi2022} A. D. Bordoloi 等, \emph{Nat. Commun.} \textbf{13}, 3820 (2022).
			\bibitem{Kale2023} S. Kale 等, \emph{Phys. Biol.} \textbf{20}, 046002 (2023).
			\bibitem{Stauber2015} J. M. Stauber 等, \emph{Langmuir} \textbf{31}, 2353 (2015).
			\bibitem{Blatter1994} G. Blatter 等, \emph{Rev. Mod. Phys.} \textbf{66}, 1125 (1994).
			\bibitem{USGSNEIC} USGS NEIC 地震目录.
			\bibitem{FishBase} FishBase POPGROWTH 表.
			\bibitem{TN1} A. V. Yakushev, 技术说明 1：反常扩散、碎裂、玻璃松弛 (2026).
			\bibitem{TN2} A. V. Yakushev, 技术说明 2：细胞表面–体积、晶体生长、超导电流 (2026).
			\bibitem{TN3} A. V. Yakushev, 技术说明 3：陨石坑、地震、液滴蒸发 (2026).
			\bibitem{TN4} A. V. Yakushev, 技术说明 4：鱼类生长、城市规模、代谢标度 (2026).
			\bibitem{TN5} A. V. Yakushev, 技术说明 5：细菌标度、晶体生长、超导 (2026).
		\end{thebibliography}
		
	\end{multicols}
\end{document}